\documentclass[10pt]{article}
\usepackage[T1]{fontenc}
\usepackage[utf8]{inputenc}
\usepackage[french]{babel}
\usepackage{amsmath,amssymb,mathrsfs,mathpazo}
\usepackage{geometry}
\geometry{a4paper,margin=2cm}
\usepackage{graphicx,array,tabularx,tabularray,colortbl,multirow,multicol,booktabs,enumitem,pifont,ulem,fancybox,diagbox,cancel}
\usepackage{pstricks,pstricks-add,pst-plot,pst-tree,pst-node,pst-text,pst-func,pst-math,pst-bspline,pst-3dplot,pst-eucl,pst-solides3d}
\let\pscalculate\undefined
\usepackage{pst-bezier}
\usepackage{tikz,pgfplots,tkz-tab}
\usetikzlibrary{arrows,arrows.meta,patterns,calc,positioning,trees,decorations.markings,shapes.geometric,shapes.misc,shapes.arrows,intersections,decorations.pathmorphing,decorations.pathreplacing,backgrounds,shadows,babel}
\pgfplotsset{compat=1.18}
\usepgfplotslibrary{fillbetween}
\usepackage[np]{numprint}
\usepackage{listings,makecell,pas-tableur,twoopt}
\let\barre\undefined
\usepackage[pstricks]{tablvar}
\definecolor{darkgreen}{RGB}{0,100,0}
\providecommand{\red}{\color{red}}
\providecommand{\blue}{\color{blue}}
\providecommand{\green}{\color{green}}
\providecommand{\black}{\color{black}}
\providecommand{\white}{\color{white}}
\providecommand{\gray}{\color{gray}}
\providecommand{\cyan}{\color{cyan}}
\providecommand{\magenta}{\color{magenta}}
\providecommand{\yellow}{\color{yellow}}
\providecommand{\violet}{\color{violet}}
\providecommand{\orange}{\color{orange}}
\providecommand{\euro}{\texteuro}
\providecommand{\EUR}{\texteuro}
\providecommand{\ieme}{\textsuperscript{e}}
\providecommand{\degres}{\ensuremath{^\circ}}
\providecommand{\degree}{\ensuremath{^\circ}}
\providecommand{\vv}[1]{\overrightarrow{#1}}
\providecommand{\Coord}[1]{#1}
\providecommand{\decosix}{}
\providecommand{\starredbullet}{\ensuremath{\bullet}}
\pagestyle{empty}
\setlength{\parindent}{0pt}
\usepackage{adjustbox}\DeclareUnicodeCharacter{2236}{:}\DeclareUnicodeCharacter{2190}{\ensuremath{\leftarrow}}
\providecommand{\R}{\mathbb{R}}
\providecommand{\N}{\mathbb{N}}
\providecommand{\Z}{\mathbb{Z}}
\providecommand{\Q}{\mathbb{Q}}
\providecommand{\D}{\mathbb{D}}
\providecommand{\C}{\mathbb{C}}
\providecommand{\e}{{\mathrm{e}}}
\providecommand{\vect}[1]{\overrightarrow{#1}}
\providecommand{\vectt}[1]{\overrightarrow{\mathrm{#1}}}
\providecommand{\vvt}[1]{\overrightarrow{\mathrm{#1}}}
\providecommand{\np}[2][]{{#2~#1}}
\providecommand{\vv}[1]{{\overrightarrow{#1}}}
\renewcommand{\d}{{\mathrm{d}}}
\providecommand{\degres}{{^\circ}}
\providecommand{\degre}{{^\circ}}
\providecommand{\blue}{\color{blue}}
\providecommand{\red}{\color{red}}
\providecommand{\enskip}{\quad}
\providecommand{\llbracket}{\lbrack\!\lbrack}
\providecommand{\rrbracket}{\rbrack\!\rbrack}
\providecommand{\euro}{{\text{€}}}
\providecommand{\no}{{\text{n°}}}
\providecommand{\interval}[2]{{[#1;#2]}}
\providecommand{\intervFO}[2]{{[#1;#2[}}
\providecommand{\dotfill}{\ldots}

\providecommand{\ds}{\displaystyle}
\providecommand{\textding}[1]{#1}
\providecommand{\Oij}{\(\mathrm{O};\vec{\imath},\vec{\jmath}\)}
\providecommand{\Oijk}{\(\mathrm{O};\vec{\imath},\vec{\jmath},\vec{k}\)}
\providecommand{\Ouv}{\(\mathrm{O};\vec{u},\vec{v}\)}
\providecommand{\intervFF}[2]{\ensuremath{[#1;#2]}}
\renewcommand{\euro}{\ifmmode\text{\texteuro}\else\texteuro\fi}
\DeclareUnicodeCharacter{00B0}{\ifmmode^{\circ}\else\textdegree\fi}
\DeclareUnicodeCharacter{20AC}{\ifmmode\text{\texteuro}\else\texteuro\fi}
% Énoncés : Éducation nationale. Transcription et corrigés : APMEP, Denis Vergès et contributeurs. Découpage et classement : CoopMaths, Sébastien Lozano, Éric Elter et contributeurs. Adaptation de la présentation : Maths Explorer.
% CC BY-SA : https://coopmaths.fr/www/about/
% https://coopmaths.fr/alea/static/bac/2023/tex/bac_2023_03_sujet2_polynesie_1_cor.tex
\begin{document}
\noindent\textit{APMEP / CoopMaths. CC BY-SA. Présentation adaptée par Maths Explorer.}\par\medskip
\bigskip

\textbf{Partie A}
 
 \begin{enumerate}
\item  	

	\hfill~
			
			\begin{tikzpicture}[xscale=1,yscale=1,baseline=(R)]
				
				\tikzstyle{fleche}=[->,>=latex,thick]
				\tikzstyle{noeud}=[fill=white,circle,inner sep=2pt]
				\tikzstyle{feuille}=[fill=white,circle,inner sep=2pt]
				\tikzstyle{etiquette}=[midway,fill=white, inner xsep=3pt, inner ysep=1.5pt]
				
				\def\DistanceInterNiveaux{3}
				\def\DistanceInterFeuilles{0.8}
				
				\def\NiveauA{(0)*\DistanceInterNiveaux}
				\def\NiveauB{(1)*\DistanceInterNiveaux}
				\def\NiveauC{(2)*\DistanceInterNiveaux}
				\def\InterFeuilles{(-1)*\DistanceInterFeuilles}
				
				\node[noeud] (R) at ({\NiveauA},{(1.5)*\InterFeuilles}) {};
				\node[noeud] (Ra) at ({\NiveauB},{(0.5)*\InterFeuilles}) {$R_n$};
				\node[feuille] (Raa) at ({\NiveauC},{(0)*\InterFeuilles}) {$R_{n+1}$};
				\node[feuille] (Rab) at ({\NiveauC},{(1)*\InterFeuilles}) {$\overline{R_{n+1}}$};
				\node[noeud] (Rb) at ({\NiveauB},{(2.5)*\InterFeuilles}) {$\overline{R_n}$};
				\node[feuille] (Rba) at ({\NiveauC},{(2)*\InterFeuilles}) {$R_{n+1}$};
				\node[feuille] (Rbb) at ({\NiveauC},{(3)*\InterFeuilles}) {$\overline{R_{n+1}}$};
				
				\draw[fleche] (R.east)--(Ra.west) node[etiquette] {$p_n$};
				\draw[fleche] (Ra.east)--(Raa.west) node[etiquette] {$0,9$};
				\draw[fleche] (Ra.east)--(Rab.west) node[etiquette] {$0,1$};
				\draw[fleche] (R.east)--(Rb.west) node[etiquette] {$1 - p_n$};
				\draw[fleche] (Rb.east)--(Rba.west) node[etiquette] {$0,3$};
				\draw[fleche] (Rb.east)--(Rbb.west) node[etiquette] {$0,7$};
			\end{tikzpicture}
	\hfill~
		

\item Les évènements $R_n$ et $\overline{R_n}$ partitionnent l'univers, la loi des probabilités totales donne :

$p_{n+1} = P(R_{n+1}) = P\Big(R_n \cap R_{n+1}\Big) + P\Big(\overline{R_n} \cap R_{n+1} \Big) = p_n \times 0,9 + (1-p_n) \times 0,3 = 0,9p_n + 0,3 - 0,3p_n $

Donc on a bien $p_{n+1} = 0,6p_n + 0,3$ : on a donc bien établi la relation de récurrence annoncée.

\item
\begin{enumerate}
\item  Établissons la relation de récurrence de la suite $(u_n)$. Soit $n$ un entier naturel :

$\begin{aligned} u_{n+1} &= p_{n+1} - 0,75 \quad \text{par définition de la suite } (u_n)\\
				&= (0,6p_n + 0,3) - 0,75 \quad \text{par la relation de récurrence de la suite } (p_n)\\
				&= 0,6p_n - 0,45\\
				&= 0,6(u_n + 0,75)-0,45\quad \text{par définition de la suite } (u_n)\\
				&=0,6u_n+0,45-0,45\\
				&= 0,6u_n \\ \end{aligned}$

				La relation de récurrence de la suite $(u_n)$ est donc bien celle d'une suite géométrique, de raison $q = 0,6$. Le premier terme de la suite est $u_0 = p_0 - 0,75 = 0,6 - 0,75 = - 0,15$.
\item On peut donc donner la forme explicite du terme général de la suite géométrique :

$\forall n \in \N, \quad u_n = u_0 \times q^n = -0,15\times 0,6^n$.

On en déduit : $u_n = p_n - 0,75 \iff p_n = u_n + 0,75 = 0,75 - 0,15\times 0,6^n.$
\[\red p_n = 0,75 - 0,15\times 0,6^n.\]
\item La raison de la suite géométrique $u$ est comprise entre $-1$ et 1, strictement, donc la suite $u$ converge vers 0.

Par limite de la somme, on en déduit que la suite $p$ converge vers $\ell=0,75$.
\item En assimilant les probabilités à des proportions,  au bout d'un certain temps,  l'athlète franchira la barre  dans 75\,\% des cas.
\end{enumerate}

\end{enumerate}

\textbf{Partie B}

\begin{enumerate}
\item Ici, nous avons :
\begin{itemize}
\item une expérience à deux issues  dont le succès est \og l'athlète  franchit la haie \fg{} a une probabilité $p = 0,75$;
\item cette expérience est répétée dix fois (la course comporte 10 haies) d'une façon assimilable à une répétition identique et indépendante , donc $n=10$ répétitions;
\item dans ce schéma de Bernoulli de paramètres $n=10$ et $p=0,75$, on note $X$ la variable aléatoire qui compte le nombre de succès (le nombre de haies franchies pendant la course).
\end{itemize}
Avec ces éléments, on peut dire que $X$ suit la loi binomiale $ \mathcal{B}\left(10 ~;~ 0,75\right) $.
\item La probabilité demandée est $P(X = 10)$. Par propriété, on a :

$\displaystyle P(X=10) = \binom{10}{10} \times 0,75^{10} \times 0,25^0 = 0,75^{10}\approx 0,056$.

La probabilité que l'athlète franchisse les dix haies est donc d'environ  0,056.

\item  La probabilité demandée est $P(X \geqslant 9)$.

Selon le modèle de calculatrice, on peut obtenir une valeur approchée de ce résultat directement, ou alors on a recours à : $P(X \geqslant 9) = 1 - P(X <9) = 1 - P(X \leqslant 8)$.

La calculatrice donne une valeur approchée au millième près qui est 0,244.
\end{enumerate}

\bigskip

\end{document}
